\subsection{The Postnikov Tower of a Cubeset}\label{ssec:postnikov}

Now we have all the techniques at hand which allow one to prove the weak structure theorem.
This theorem gives insight in the structure of concrete fibrant cubesets.
It appears as \cite[Theorem 3.2.19]{Candela2017} and \cite[Corollary 7.20]{Gutman2020}
and in both cases its proof covers large parts of the papers.

\begin{theorem}[Weak Structure Theorem]\label{thm:weak-structure}
Let $X$ be a concrete ergodic fibrant cubeset.
Then there exists a tower of surjective fibrations
\begin{center}
\begin{tikzcd}
	X & \cdots & {\tau_2X} & {\tau_1X} & {\tau_0X}
	\arrow[from=1-1, to=1-2]
	\arrow[from=1-2, to=1-3]
	\arrow[from=1-3, to=1-4]
	\arrow[from=1-4, to=1-5]
\end{tikzcd}
\end{center}
such that $\tau_n X$ is a concrete fibrant $n$-step cubeset
and $\tau_n X$ is a principal $\cD_n(\pi_n(X,x_0))$-bundle over $\tau_{n-1} X$.
This tower is natural in $X$.
\end{theorem}

By principal $\cD_n(\pi_n(X,x_0))$-bundle we mean principal bundle in the wide sense,
i.e., without the assumption of local triviality.
In general, given a group object $G$ in a category $\mcC$,
a \emph{$G$-torsor} is an object $X$ together with a $G$-action $\alpha\colon G\times X\to X$
that is \emph{free} and \emph{transitive}, meaning that the \emph{shear map}
\[
    \alpha\times\pi_2\colon G\times X\to X\times X,
\]
which arises from the action $\alpha$ and the projection $\pi_2$ onto $X$, is an isomorphism.
Informally, surjectivity corresponds to transitivity and injectivity is freeness.
Now a \emph{principal $G$-bundle over $X$} is a relative $G$-torsor over a base space $X$,
i.e., a $G$-torsor in the slice category $\mcC_{/X}$.
Spelled out, this means that it is a map $p\colon P\to X$ together with a $G$-action
$\alpha\colon G\times P\to P$ over $X$ such that the shear map
\[
    \alpha\times\pi_2\colon G\times P\to P\times_X P
\]
is an isomorphism and the diagram
\begin{center}
\begin{tikzcd}
	{G\times P} & P & X
	\arrow["\alpha", shift left, from=1-1, to=1-2]
	\arrow["{\pi_2}"', shift right, from=1-1, to=1-2]
	\arrow["p", from=1-2, to=1-3]
\end{tikzcd}
\end{center}
is a coequalizer diagram.
Note that if $\mcC$ is a (Grothendieck) topos,
these conditions are equivalent to asserting that $p\colon P\to X$ is an epimorphism
and that the shear map is an isomorphism, by Giraud's axioms, see \cite[Appendix]{MacLane1994}.
And we say that $p$ is \emph{locally trivial} if there exists a cover $U\to X$ such that
there is a pullback diagram of the form
\begin{center}
\begin{tikzcd}
	{G\times U} & P \\
	U & X
	\arrow[from=1-1, to=1-2]
	\arrow["{\pi_2}"', from=1-1, to=2-1]
	\arrow["p", from=1-2, to=2-2]
	\arrow[from=2-1, to=2-2]
\end{tikzcd}.
\end{center}
If $\mcC$ is an $\infty$-topos, the situation becomes substantially easier
since colimits are better behaved.
Namely, it suffices to require that $X = P/G$
and then the shear map automatically is an equivalence,
and furthermore the principal bundle is locally trivial, see \cite{Nikolaus2014}.

\begin{remark}
If $X$ is not necessarily ergodic, then we can decompose $X$ into its ergodic components,
see Proposition~\ref{prop:ergodic-decomp}.
Applying the weak structure theorem to each of these components and noting the truncation
$\tau_n$ preserves the coproduct, see Proposition~\ref{prop:n-trunc-closure},
we obtain the weak structure theorem for non-ergodic cubesets.
However, then the fibration $\tau_nX\to\tau_{n-1}X$ is not necessarily a principal bundle for a group anymore,
but rather a principal groupoid bundle for the groupoid $\pi_n(X)$.

Since this variant is notationally more elaborate we won't formulate it here.
But we note that morally, $\pi_n(X)$ shouldn't just be a groupoid, but an $n$-groupoid,
and the whole $X$ an $\infty$-groupoid (= anima).
This view on $X$ is taken in \cite{Bihlmaier2026a}.
\end{remark}

To establish the weak structure theorem take a concrete fibrant ergodic cubeset $X$.
Let us look at the unit $\eta\colon X\to\tau_{n-1}X$
where $X$ is $n$-step.
Since $\eta$ is a surjective fibration we obtain a fiber sequence for every $x_0\in X$:
\begin{center}
\begin{tikzcd}
	{(E,x_0)} & {(X,x_0)} & {(\tau_{n-1}X,x_0)}
	\arrow["e", from=1-1, to=1-2]
	\arrow["\eta", from=1-2, to=1-3]
\end{tikzcd}.
\end{center}

\begin{proposition}\label{prop:fiber-ergodic-trunc}
The cubeset $E$ is a concrete fibrant cubeset that is both $n$-ergodic and $n$-step.
Moreover, $e_\ast\colon\pi_n(E,x_0)\to\pi_n(X,x_0)$ is an isomorphism.
\end{proposition}

\begin{proof}
Subobjects of concrete cubesets are concrete, so $E$ is concrete.

Now for $k>n$, given a $k$-corner $\lambda$ in $E$, one can complete it in $\tau_nX$ to a $k$-cube $c$ (since $\tau_nX$ is fibrant).
Under $\eta$, the $k$-corner $e\circ\lambda$ of $c$ is mapped to $x_0$.
By unique corner completion in $\tau_{n-1}X$ we must have that $\eta\circ c \equiv x_0$ and so $c$ factors through $e$, i.e., $c\in E$.
\begin{center}
\begin{tikzcd}
	{{\bbcor}^k} & E & {X} & {\tau_{n-1}X} \\
	{{\bbox}^k} & \ast & \ast & \ast
	\arrow["\lambda", from=1-1, to=1-2]
	\arrow[hook', from=1-1, to=2-1]
	\arrow["e", from=1-2, to=1-3]
	\arrow[from=1-2, to=2-2]
	\arrow["\eta", from=1-3, to=1-4]
	\arrow[from=1-3, to=2-3]
	\arrow[from=1-4, to=2-4]
	\arrow[dashed, from=2-1, to=1-2]
	\arrow["c"{description, pos=0.3}, dashed, from=2-1, to=1-3]
	\arrow[from=2-1, to=2-2]
	\arrow["{=}"{description}, draw=none, from=2-2, to=2-3]
	\arrow["{=}"{description}, draw=none, from=2-3, to=2-4]
\end{tikzcd}
\end{center}
Uniqueness of the completion in $E$ follows from uniqueness in $X$.

Given $x,y\in E(0)$ we have that $\eta(x) = \eta(y) = x_0$ in $\tau_{n-1}X$
and so $x$ and $y$ are ($n-1$)-equivalent in $X$, i.e., $x\sim_{n-1} y$.
Using universal replacement~\ref{prop:univ-replacement} we see that
we can arbitrarily interchange $x$ and $y$ in every $k$-cube of $X$ for $k\leq n$.
Therefore, every map $\{0,1\}^k\to E(0)$ is in $E(k)$ and so $E(k) \simeq \prod_{0\to k} E(0)$ for $k\leq n$,
hence $E$ is $n$-ergodic.
From the previous two items it also follows that $E$ is fibrant.

The last assertion follows from the fact that $X$ is $n$-step,
hence $\pi_n(X,x_0) = \{c\in X(n)\mid c\rvert_{{\bbcor}^n}\equiv x_0\}$.
Given such a $c\in\pi_n(X,x_0)$ we have that $\tau_{n-1}(c)\equiv x_0$ and thus $c\in E(n)$.
\end{proof}

Now it remains to show that $E\simeq\cD_n(A)$ for some abelian group $A$.
Which abelian group is available? Of course the structure group $\pi_n(E,x_0) = \pi_n(X,x_0)$.

\begin{proposition}\label{prop:fiber-torsor}
Let $E$ be an $n$-ergodic $n$-step concrete cubeset.
Let $x_0\in E$ and $A = \pi_n(E,x_0)$.
Then there is an action $\alpha\colon\cD_n(A)\times E\to E$ such that $E$ is an $\cD_n(A)$-torsor.
\end{proposition}

\begin{proof}
We have that $A = \pi_n(E,x_0) = \{ c\in E(n) \mid c\rvert_{{\bbcor}^n} \equiv x_0 \}$ because $E$ is $n$-step.
Since $E$ is $n$-ergodic and concrete, evaluation at the uppermost vertex induces a bijection $A\simeq E(0)$.
This allows us to define an action
\[
    \alpha_0\colon A\times E(0)\to E(0)
\]
which is just given by applying the bijection $E(0)\simeq A$ and then acting by left multiplication on $A$,
and afterwards forgetting again to $E(0)$.
For $k\leq n$ this defines by pushforward of $\alpha_0$ an action
\[
    \alpha_k\colon\cD_n(A)(k)\times E(k) = \hom_\set(\{0,1\}^k,A\times E(0)) \to \hom_\set(\{0,1\}^k,E(0)) = E(k)
\]
natural in $k$ (which is just the pointwise action).
For $k\geq n+1$ recall that we can write
\[
    {\bbcor}^k = \varinjlim_{j\in\cJ}{\bbox}^j
\]
where $\cJ$ runs over all inert inclusions ${\bbox}^j\to{\bbox}^k$ for $j<k$.
Suppose we have defined the action $\alpha_m\colon\cD_n(A)(m)\times E(m)\to E(m)$ for all $m<k$ natural in $m$.
This yields a map
\[
    \varprojlim_{j\in\cJ}\cD_n(A)(j)\times E(j) \to \varprojlim_{j\in\cJ} E(j)
\]
which then after using the isomorphism
\begin{align*}
    \cD_n(A)(k)\times E(k)
    &= \hom({\bbox}^k,\cD_n(A)\times E) \\
    &\simeq \hom({\bbcor}^k,\cD_n(A)\times E) \\
    &= \varprojlim_{j\in\cJ}\hom({\bbox}^j,\cD_n(A)\times E) \\
    &= \varprojlim_{j\in\cJ}\cD_n(A)(j)\times E(j)
\end{align*}
determines an action $\alpha_k\colon\cD_n(A)(k)\times E(k)\to E(k)$
(it is an action since limits commute with limits and, in particular, products) natural in $m\leq k$.
Putting every dimension together, we obtain an action
\[
    \alpha\colon\cD_n(A)\times E\to E.
\]
We go on to show that the shear map is an isomorphism.
This is clear for $k = 0$ using the bijection $A\simeq E(0)$.
But then inductively it follows for all $k\in\N$
because pushforward and taking limits is functorial, and so preserves isomorphisms.
Since $E$ is non-empty, we infer that also $E\simeq\cD_n(A)$ as cubesets by general nonsense:
every non-empty $G$-torsor is (non-canonically) isomorphic to $G$.
\end{proof}

The action $\alpha$ on $E$ does, a priori, depend on the choice of basepoint $x_0$.
Since $\pi_n(E,x_0)\simeq\pi_n(E,y_0)$ for any choice of basepoints $x_0,y_0\in E$,
we obtain an action of $\pi_n(E,y_0)$ on $E$.
Moreover, the action is compatible with this isomorphism in the following sense.
Recall from Proposition~\ref{prop:structure-monoid-independence-base-point}
that if $x_0,y_0\in E$ are two points then there is a canonical group isomorphism
$h_{(x_0,y_0)}\colon\pi_n(E,x_0)\to\pi_n(E,y_0)$.

\begin{lemma}\label{lem:action-independence-basepoint}
Let $E$ be an $n$-ergodic $n$-step concrete cubeset.
Let $x_0,y_0\in E$ and $A = \pi_n(E,x_0)$, $B = \pi_n(E,y_0)$.
Then the diagram
\begin{center}
\begin{tikzcd}
	{\cD_n(A)\times E} && {\cD_n(B)\times E} \\
	& E
	\arrow["{h\times\id}", from=1-1, to=1-3]
	\arrow["\alpha"', from=1-1, to=2-2]
	\arrow["\beta", from=1-3, to=2-2]
\end{tikzcd}
\end{center}
is commutative where $h = \cD_n(h_{(x_0,y_0)})$
and $\alpha$, $\beta$ are defined as in (the proof of) Proposition~\ref{prop:fiber-torsor}.
\end{lemma}

\begin{proof}
Since $E$ and $\cD_n(A)$ are concrete it suffices to see commutativity on level $k=0$
which means that
\[
    \beta_0\circ(h_{(x_0,y_0)}\times\id_0) = \alpha_0.
\]
Let $x\in E(0)$ and $a\in A$.
Write $x_{y_0}\in E(n)$ for the corresponding element in $\pi_n(E,y_0)$,
and analogously $x_{x_0}\in\pi_n(E,x_0)$ for $x_0$.
By definition,
\[
    \beta_0\circ(h_{(x_0,y_0)}\times\id_0)(a,x) = (h_{(x_0,y_0)}(a)*x_{y_0})(1^n)
    \quad\text{and}\quad
    \alpha_0(a,x) = (a*x_{x_0})(1^n).
\]
We compute that
\begin{align*}
       (a*x_{x_0})(1^n)
    &= (h_{(x_0,x)}(a))(1^n) \\
    &= (h_{(y_0,x)}(h_{(x,y_0)}(h_{(x_0,x)}(a))))(1^n) \\
    &= (h_{(x,y_0)}(h_{(x_0,x)}(a))*x_{y_0})(1^n) \\
    &= (h_{(x_0,y_0)}(a)*x_{y_0})(1^n),
\end{align*}
where we have used Lemma~\ref{lem:homotopy-action-well-defined}
in the first and third equality and Lemma~\ref{lem:homotopy-action-well-defined-global}
in the second and fourth equality.
\end{proof}

In the last step we show that the map $X\to\tau_{n-1}X$ is a principal $\cD_n(A)$-bundle
for $X$ $n$-step and $A = \pi_n(X,x_0)$.

\begin{proposition}\label{prop:unit-principal-bundle}
Let $X$ be an ergodic concrete fibrant $n$-step cubeset, $x_0\in X(0)$ and $A = \pi_n(X,x_0)$.
Then the unit $\eta\colon X\to\tau_{n-1}X$ is a principal $\cD_n(A)$-bundle.
\end{proposition}

\begin{proof}
We begin by defining an action $\alpha_0\colon A\times X(0)\to X(0)$.
Given $x\in X(0)$, by ergodicity we have that $(x_0,x)\in X(1)$.
Recall the monoid isomorphism $h_{(x_0,x)}\colon\pi_n(X,x_0)\to\pi_n(X,x)$ from Proposition~\ref{prop:structure-monoid-independence-base-point}.
Then define
\[
    \alpha_0((a,x)) = h_{(x_0,x)}(a)(1^n)\in X(0).
\]
This defines an action of $A$ on $X(0)$:
since $h_{(x_0,x)}$ is a monoid homomorphism we have that
\[
    \alpha_0((e_{x_0},x)) = h_{(x_0,x)}(e_{x_0})(1^n) = e_x(1^n) = x.
\]
Associativity follows from Lemma~\ref{lem:homotopy-action-well-defined}:
given $a,b\in A$ and $x\in X(0)$ we obtain for $z_0 = \alpha_0(b,x) = h_{(x_0,x)}(b)(1^n)$ that
$z_0\sim_{n-1} x$ and thus
\begin{align*}
       \alpha_0(a,\alpha_0(b,x))
    &= h_{(x_0,z_0)}(a)(1^n) \\
    &= h_{(x,z_0)}(h_{(x_0,x)}(a))(1^n) \\
    &= (h_{(x_0,x)}(a)*h_{(x_0,x)}(b))(1^n) \\
    &= h_{(x_0,x)}(a*b)(1^n)
    = \alpha_0(a*b,x)
\end{align*}
where we have used Lemma~\ref{lem:homotopy-action-well-defined-global} in the second equality.
Moreover, the action leaves the fibers of $\eta$ over $Y = \tau_{n-1}X$ invariant:
By definition, the $n$-cube $h_{(x_0,x)}(a)$ satisfies $h_{(x_0,x)}(a)(\omega) = x$ for all $\omega\neq 1$
and hence $\alpha_0(a,x)\sim_{n-1} x$, which in turn means that $\eta_0(\alpha_0(a,x)) = \eta_0(x)$.

Since $i_\ast\colon\set\to\cSet$ preserves finite products,
we obtain an action $i_\ast(\alpha_0)\colon i_\ast(A\times X(0))\to i_\ast X(0)$ of (concrete) cubesets.
Since $X$ and $\cD_n(A)$ are concrete it suffices to show that we obtain a factorization
\begin{center}
\begin{tikzcd}
	{\cD_n(A)\times X} & {i_\ast(A\times X(0))} \\
	X & {i_\ast X(0)}
	\arrow[hook, from=1-1, to=1-2]
	\arrow["\alpha"', dashed, from=1-1, to=2-1]
	\arrow["{i_\ast(\alpha_0)}", from=1-2, to=2-2]
	\arrow[hook, from=2-1, to=2-2]
\end{tikzcd}.
\end{center}
Then clearly the action has to arise pointwise,
so we are left to show that this is well-defined.
For $k\leq n$ we use the fact that $\alpha_0$ leaves the fibers of $\eta_0$ invariant
and so we can use universal replacement for $\sim_{n-1}$ to obtain the desired result.
For $k\geq n+1$ we use unique corner completion in $X$ to define the action
pointwise on the corner and then complete the cube,
analogously to the argument in the proof of Proposition~\ref{prop:fiber-torsor}.
This also directly shows that the action is invariant over $\eta$ by checking it pointwise
which means that $\eta\circ\alpha = p_Y\circ(\id\times\eta)$.

Moreover, the action has the following property, analogous to Lemma~\ref{lem:action-independence-basepoint}:
if $x_0,y_0\in X$ and $A = \pi_n(X,x_0)$, $B = \pi_n(X,y_0)$ then the diagram
\begin{center}
\begin{tikzcd}
	{\cD_n(A)\times X} && {\cD_n(B)\times X} \\
	& X
	\arrow["{h\times\id}", from=1-1, to=1-3]
	\arrow["\alpha"', from=1-1, to=2-2]
	\arrow["\beta", from=1-3, to=2-2]
\end{tikzcd}
\end{center}
is commutative where $h = \cD_n(h_{(x_0,y_0)})$.
This follows by checking it on points (since everything is concrete) from Lemma~\ref{lem:homotopy-action-well-defined-global}.

It remains to show that each fiber $E$ is an $\cD_n(A)$-torsor:
By invariance over $\eta$ there is an induced action $\cD_n(A)\times E\to E$.
First assume that $x_0\in E$.
Then from Lemma~\ref{lem:action-independence-basepoint} we see that
the restricted action agrees with the action defined in Proposition~\ref{prop:fiber-torsor}.
In particular, $E$ is an $\cD_n(A)$-torsor.
Otherwise, shift the action by $h$ to a basepoint in $E$ using the previous observation
and then argue analogously.
\end{proof}

\begin{proof}[Proof of Theorem~\ref{thm:weak-structure}]
The existence of the tower of fibrations is clear:
since adjunctions and hence units compose, each unit $X\to\tau_n X$ and $\tau_n X\to\tau_{n-1} X = \tau_{n-1}\tau_n X$
are surjective fibrations by Proposition~\ref{prop:trunc_concfib}.
The same proposition also asserts that $\tau_n X$ is a concrete fibrant $n$-step cubeset.
Naturality in $X$ follows from functoriality of $\tau_n$.

It remains to show that $\tau_n X\to\tau_{n-1} X$ is a principal $\cD_n(\pi_n(X,x_0))$-bundle.
First note that since $X$, and hence $\tau_n X$, is ergodic, $A = \pi_n(X,x_0)$ doesn't depend on $x_0$
by Proposition~\ref{prop:structure-monoid-independence-base-point}.
Moreover, $\pi_n(X,x_0) = \pi_n(\tau_n X,x_0)$ by Proposition~\ref{prop:trunc-structure-monoid}
and so Proposition~\ref{prop:unit-principal-bundle} gives the claim.
\end{proof}

Note that, different from the situation of simplicial sets and their Postnikov towers,
this construction of the tower does not converge as it identifies too much information.

\begin{example}
Let $S$ be a nontrivial set and let $X = i_\ast S$, i.e., $X(n) = S^{2^n}$.
Then $\tau_n X = \ast$ for every $n\in\N_0$ and so the canonical map $X\to\varprojlim_n\tau_nX = \ast$ is not an isomorphism.
\end{example}

However, it is a weak equivalence.

\begin{definition}
A map $f\colon X\to Y$ between cubesets is called a \emph{weak equivalence}
if it induces a bijection $f_\ast\colon\pi_0(X)\to\pi_0(Y)$ and
monoid isomorphisms $f_\ast\colon\pi_n(X,x_0)\to\pi_n(Y,f(x_0))$ for all $x_0\in X$ and $n\in\N$.
\end{definition}

\begin{remark}[Model Category Structures]
It is a very canonical question to ask whether there exists the structure of a model category on the category of cubesets that has this notion of weak equivalences, and/or the notion of fibration from Definition~\ref{def:nil-fibration} as fibrations.
Whilst we are not able to fully answer this question, let us remark on some obstructions,
indicating that for the present choice of fibrations from Definition~\ref{def:nil-fibration}
the answer might be negative.
\begin{enumerate}[(i)]
    \item There is no model structure on $\cSet$ with the present choice of fibrations with cofibrations being the class of monomorphisms.
		This follows from Cisinski's \emph{model structures ex nihilo} machinery in \cite{Cisinski2025}: The subobject classifier $\Omega$
		has the problem that the trivial cofibrations $\ell(r(\Lambda))$ do not consist of $\Omega$-anodyne extensions,
since in general $\ihom(\Omega, X)\to X\times_Y \ihom(\Omega, Y)$ is not surjective.
    \item We conjecture that the structure groups of the fibrant replacement of the point, $\ast_{\mathrm{fib}}$, are non-trivial.
		Thus, there should exist a trivial cofibration, e.g., $\ast\to \ast_{\mathrm{fib}}$,
		that does \emph{not} induce isomorphisms on structure groups.
		This shows that the present choice of fibrations together with the canonical choice of weak equivalences as those maps that induce isomorphisms on all structure groups
		does not induce a model structure.
	\end{enumerate}

	We were not able to disprove if using either normal monomorphisms as cofibrations, or using
	\[
    \Gamma = \{{\bbox}^n\sqcup_{{\bbcor}^n}{\bbox}^n\to{\bbox}^{n+1}\sqcup_{{\bbcor}^{n+1}}{\bbox}^{n+1} : n\in\N_0\cup\{-1\}\}
\]
	as generating cofibrations (see Remark~\ref{rem:nil-cofibrations}), yields a (useful) model structure on $\cSet$.
\end{remark}
